% ===================================================================
% Thesis summary sheet and pointers (School of EE&T form, "Thesis Project
% summary sheet S2-16"). Rules from the form's guidelines:
%   - 12 point fonts throughout (the document is 12pt; nothing here is smaller)
%   - the summary sheet must not extend beyond one page
%   - no new information that is not in the report
%   - contributions: most important first, at most one per line
%   - pointers: page numbers on the left; contributions in the same order as D;
%     literature: reference identifier on the right, citing pages on the left
%     with the page of most impact leftmost
% Page numbers are \pageref, so they follow the document. Sign and date by hand.
% ===================================================================
\pagestyle{empty}
\begingroup
\linespread{1.05}\selectfont
\setlength{\parindent}{0pt}
\newcommand{\sumbox}[1]{\noindent\fbox{\parbox{\dimexpr\textwidth-2\fboxsep-2\fboxrule\relax}{#1}}\par}
\newcommand{\ptr}[2]{#1 & #2\\ \hline}

\textbf{Topic Title:}\quad Comparison and Implementation of Data Compression Methods for SoC-based Radar Systems\\[2mm]
\textbf{Student Name:}\quad Jordan Mohindra \hfill \textbf{Student ID:} z5418455\\[3mm]

A.\quad\textbf{Problem statement}\\
\sumbox{In automotive radar, high-resolution cascaded MIMO sensors produce more
raw data than a vehicle's network and central processor can comfortably absorb.
Compressing at the sensor saves link bandwidth, but the compressor must be
lossless (anything it discards is lost before target detection), keep up with
the sensor, and be small enough for sensor hardware. Radar compression methods
are rarely compared on the same real multi-transmitter data, and the lossless
baseline was evaluated on a single-transmitter radar with no decoder.}
\vspace{2mm}

B.\quad\textbf{Objective}\\
\sumbox{Compare compression methods side by side on real, public,
multi-transmitter radar data; select the best lossless methods; and implement
their encoders and decoders on an FPGA, verified and measured on a board.}
\vspace{2mm}

C.\quad\textbf{My solution}\\
\sumbox{Eight methods were run in MATLAB on 50 ColoRadar frames (Texas
Instruments cascade, 12 transmitters, 16 receivers) through one chain: range
FFT, fixed 16-bit (FX16) scaling, compression, Doppler FFT, non-coherent
integration and peak detection, against a double-precision reference. The two
lossless predictive coders, DRHE and a second-order slow-time LPC, predict from
the same transmitter (12 ramps back) and share an $S_4$/Huffman coder whose
dictionary was re-derived from five frames of their prediction errors. Both,
with matching decoders and a range-FFT front end, were written in C++ for
Vitis HLS, built with Vivado into a DMA test system on an AMD KCU116 board at
\SI{100}{\mega\hertz}, and verified bit-exact from MATLAB through C-simulation
and RTL co-simulation to the board.}
\vspace{2mm}

D.\quad\textbf{Contributions} (at most one per line, most important first)\\
\begin{tabular}{|p{\dimexpr\textwidth-2\tabcolsep-2\arrayrulewidth\relax}|}\hline
Bit-exact lossless radar compression on an FPGA: 4.22$\times$ (DRHE), 4.40$\times$ (LPC)\\ \hline
Same-transmitter prediction lag for time-division MIMO radar\\ \hline
Comparison of eight compression methods on real cascade radar data\\ \hline
Breakdown of the compression ratio; Huffman dictionaries matched to the errors\\ \hline
Encoders, decoders and range FFT measured on the FPGA board\\ \hline
Reconstruction of, and comparison with, the baseline design\\ \hline
\end{tabular}
\vspace{3mm}

E.\quad\textbf{Suggestions for future work}\\
\begin{tabular}{|p{\dimexpr\textwidth-2\tabcolsep-2\arrayrulewidth\relax}|}\hline
Fixed-point DRHE datapath to cut logic and DSP use\\ \hline
Streaming (backward-adaptive) LPC without a frame buffer\\ \hline
All sixteen receive channels and power measurement\\ \hline
\end{tabular}
\vspace{4mm}

While I may have benefited from discussion with other people, I certify that
this report is entirely my own work, except where appropriately documented
acknowledgements are included.\\[9mm]
Signature: \rule{6cm}{0.4pt} \hfill Date: \rule{1cm}{0.4pt}\,/\,\rule{1cm}{0.4pt}\,/\,\rule{1.5cm}{0.4pt}

\newpage
\enlargethispage{2\baselineskip}
\linespread{1.0}\selectfont
{\centering\textbf{\large Pointers}\par}\vspace{2mm}
List relevant page numbers in the column on the left.  Be precise and
selective: Don't list all pages of your report!\\[1mm]
\newcommand{\pcol}{2.9cm}
\newcommand{\ptab}[2]{#1\\[0.5mm]\begin{tabular}{|p{\pcol}|p{\dimexpr\textwidth-\pcol-4\tabcolsep-3\arrayrulewidth\relax}|}\hline #2\end{tabular}\\[0.6mm]}

\begin{tabular}{|p{\pcol}|p{\dimexpr\textwidth-\pcol-4\tabcolsep-3\arrayrulewidth\relax}|}\hline
\ptr{\pageref{sec:problem}}{Problem Statement}
\ptr{\pageref{sec:objectives}}{Objective}
\end{tabular}\\[0.6mm]

\ptab{\textbf{Theory} (up to 5 most relevant ideas)}{%
\ptr{\pageref{sec:tdm}}{Time-division MIMO radar and the data cube}
\ptr{\pageref{sec:compression}}{Lossless compression and entropy coding}
\ptr{\pageref{sec:fx16}}{Fixed-point (FX16) quantisation}
\ptr{\pageref{sec:fpga}}{FPGAs and high-level synthesis}
\ptr{\pageref{ch:kiem}}{The baseline: Kiem's DRHE design}}

\ptab{\textbf{Method of solution} (up to 5 most relevant points)}{%
\ptr{\pageref{sec:candidates}}{Eight candidate methods}
\ptr{\pageref{sec:entropycoder}}{Shared $S_4$/Huffman entropy coder}
\ptr{\pageref{sec:predaxis}}{Prediction on a TDM-MIMO radar}
\ptr{\pageref{sec:ladder}}{Verification ladder, MATLAB to board}
\ptr{\pageref{ch:hardware}}{Hardware implementation}}

\ptab{\textbf{Contributions} (most important first)}{%
\ptr{\pageref{sec:verifres}}{Bit-exact lossless radar compression on an FPGA}
\ptr{\pageref{sec:predaxis}, \pageref{sec:predaxisres}}{Same-transmitter prediction lag for TDM-MIMO radar}
\ptr{\pageref{sec:stage1}}{Comparison of eight compression methods}
\ptr{\pageref{sec:decomp}, \pageref{sec:dictresult}}{Breakdown of the ratio; matched Huffman dictionaries}
\ptr{\pageref{sec:boardtiming}, \pageref{sec:feresults}}{Encoders, decoders and range FFT measured on the board}
\ptr{\pageref{ch:hwinterpret}, \pageref{ch:vskiem}}{Reconstruction of, and comparison with, the baseline}}

\ptab{\textbf{My work}}{%
\ptr{\pageref{sec:drhe}--\pageref{sec:dict}, \pageref{ch:hardware}--\pageref{sec:system}}{System block diagrams/algorithms/equations solved}
\ptr{\pageref{sec:metrics}--\pageref{sec:detectionmetrics}}{Description of assessment criteria used}
\ptr{\pageref{sec:stages}, \pageref{sec:ladder}, \pageref{sec:system}}{Description of procedure (e.g.\ for experiments)}}

\ptab{Results}{%
\ptr{\pageref{sec:headtohead}, \pageref{sec:stage1}}{Succinct presentation of results}
\ptr{\pageref{sec:decomp}, \pageref{ch:hwinterpret}, \pageref{sec:ratiomeaning}}{Analysis}
\ptr{\pageref{sec:choice}}{Significance of results}}

\ptab{Conclusion}{%
\ptr{\pageref{ch:conclusion}}{Statement of whether the outcomes met the objectives}
\ptr{\pageref{sec:futurework}}{Suggestions for future research}}

\ptab{\textbf{Literature}: (up to 5 most important references)}{%
\ptr{\pageref{ch:kiem}--\pageref{ch:kiem:end}, \pageref{ch:hwinterpret}--\pageref{ch:vskiem:end}}{\cite{kiem2025} L.~Kiem, master's thesis, 2025 (the baseline design)}
\ptr{\pageref{lit:kramer:bg}, \pageref{lit:kramer:intro}}{\cite{kramer2022} A.~Kramer \emph{et al.}, ColoRadar dataset, 2022}
\ptr{\pageref{sec:drhe}, \pageref{sec:candidates}, \pageref{sec:sota}}{\cite{li2014} L.~Li, MSc thesis, 2014 (origin of DRHE)}
\ptr{\pageref{sec:kr_lpcraw}, \pageref{sec:lpc_notkiem}, \pageref{sec:kc_dse}}{\cite{meucci2022} G.~Meucci and F.~Mancuso, LPC-Huffman, 2022}
\ptr{\pageref{lit:huffman:dict}, \pageref{lit:huffman:bg}}{\cite{huffman1952} D.~A.~Huffman, minimum-redundancy codes, 1952}}
\endgroup
\clearpage
\pagestyle{plain}
\thispagestyle{empty}
